% !TeX spellcheck = en_US
% !TeX encoding = UTF-8
\chapter*{Abstract}
\markboth{Abstract}{Abstract}

This thesis builds and evaluates a weakly supervised pipeline that, for each
UCF-Crime video, decides whether it is anomalous, which of three
project-defined services (Police, Medical, Fire) it involves, and where the
event occurs. The pipeline combines VideoMAE with low-rank adaptation (LoRA),
frozen clip features, a normal-prototype distance, and an attention multiple
instance learning model with a binary head and a three-channel response head.
Each output is compared with a trivial reference on the same videos, and
Stage-2 training is repeated with three seeds over a fixed feature store.

Video-level binary detection was stable: on fixed 16-clip test bags, the area
under the receiver operating characteristic curve was $0.940\pm0.002$ (mean
$\pm$ standard deviation over the three Stage-2 seeds), a video-level value
that is not comparable with frame-level UCF-Crime results. Response selection
was weak. Without gating the head predicted Police for every video, and after
binary gating it exceeded a rule that predicts every service for each detected
anomaly by only 0.015 in Macro-F1 (95\% video bootstrap interval 0.008 to
0.023). A validation-only extension that added normal bags to the response loss
removed the universal Police output, but its gated margins stayed within the
seed spread of the control. On a dense protocol covering 99.843\% of test-video
frames, attention reached a temporal average precision of $0.327\pm0.014$ and
an interval overlap of $0.186\pm0.017$, above trivial references but far from
precise localization.

The contribution is an offline implementation and a matched evaluation that
keeps anomaly detection, response prediction, and temporal grounding separate.
The response targets are deterministic transformations of dataset categories
rather than observed service decisions, and the computer-aided dispatch
(CAD)-style output has no calibration, safety, or operational validation.
